\chapter{Web3 技术纵深：账本、共识、智能合约与市场原语}

\section{分布式账本与密码学基础}

\textbf{分布式账本技术（DLT）}指多节点共同维护、通过密码学与协议规则达成一致的数据结构，不必依赖单一中心化记账方。\textbf{公链}（如比特币、以太坊及众多 Layer~1）向任意满足规则的参与者开放验证与交易广播；\textbf{联盟链/私有链}则限制节点集合，更贴近机构间清算与供应链场景。密码学工具包括：哈希链保证篡改可检测、非对称签名保证授权与不可否认、默克尔树支撑轻客户端验证、零知识证明（ZKP）在隐私与扩容方案中日益重要。理解这些构件有助于区分「真正链上结算」与仅借用区块链品牌营销的中心化数据库。

\section{共识、结算与最终性}

\textbf{工作量证明（PoW）}通过算力竞争出块，能源与安全假设明确；\textbf{权益证明（PoS）}以质押与罚没（slashing）约束验证者行为，资本效率与委托结构不同。\textbf{最终性（finality）}指交易被撤销的概率足够低的时间点——不同链的确认数与重组风险不同，跨链与交易所入账规则必须与之对齐。\textbf{Layer~2（Rollup 等）}将执行从主链外移，通过欺诈证明或有效性证明（ZK-Rollup）向 L1 提交状态承诺，在吞吐与费用上折中，但引入排序器（sequencer）中心化与桥接风险等新攻击面。

\section{账户模型、钱包与托管}

以太坊等采用\textbf{账户模型}（EOA 合约账户）；比特币为\textbf{UTXO} 模型。钱包按密钥控制方式分为\textbf{自托管}（用户持有私钥或助记词）与\textbf{托管钱包}（交易所/银行代持）。机构场景下常见 \textbf{MPC}（多方计算）、\textbf{HSM} 与策略引擎组合，以满足职责分离与交易审批。链上地址不等于实名身份；合规上的「谁最终控制资产」须结合托管协议与链下 KYC 映射。

\section{智能合约与可组合性}

\textbf{智能合约}是在链上虚拟机中执行的确定性程序；升级机制（代理合约、多签治理）与\textbf{权限角色}（owner、admin、pause）直接决定资产安全。以太坊生态的\textbf{可组合性}允许协议在无许可情况下相互调用，加速创新也放大\textbf{级联清算}与漏洞传染风险。常见攻击类型包括：重入、预言机操纵、闪电贷辅助套利、跨链桥逻辑错误、治理劫持等——技术尽调须阅读审计报告与链上权限而不仅看白皮书。

\section{Token 标准与资产表示}

\textbf{同质化代币}（如 ERC-20）适合份额与积分式权利；\textbf{NFT}（如 ERC-721/1155）适合唯一凭证与 RWA（实物资产）映射中的编号权利。\textbf{包装资产}（如封装 BTC）依赖托管方与铸造/销毁规则，信用风险回到链下对手方。\textbf{治理代币}的经济与法律属性高度情境化：可能仅投票工具，也可能因募资方式被认定为证券（见下一章）。

\section{DeFi 核心原语（金融市场视角）}

\textbf{自动做市商（AMM）}用曲线函数替代订单簿，流动性由 LP 提供，承担无常损失与智能合约风险。\textbf{超额抵押借贷}（如稳定币铸造、循环杠杆）依赖清算机器人与预言机价格；参数（LTV、清算罚金）决定系统性压力下的级联。\textbf{稳定币机制}包括法币储备型、加密抵押型、算法型（历史波动极大）；各自脱锚与储备透明度是监管焦点。\textbf{预言机}将链外价格与事件喂入合约，是操纵与 MEV（最大可提取价值）争夺的关键环节——排序器、区块构建者与搜索者之间的利益结构，使「链上公平」成为工程与经济学交叉问题。

\section{跨链桥与互操作性}

桥将资产或消息在链间转移，技术路径包括锁定-铸造、流动性网络、轻客户端验证等。\textbf{桥}历史上是高损事件集中区：合约漏洞、多签私钥泄露、经济攻击（虚假消息）均可导致巨额盗取。机构参与跨链头寸须单独进行风险限额与应急演练。

\begin{figure}[htbp]
\centering
\begin{tikzpicture}[font=\footnotesize, node distance=0.55cm]
  \node[reg node=violet, minimum width=6.8cm] (app) {应用层：钱包、聚合器、RWA 门户、游戏};
  \node[reg node=blue, below=of app, minimum width=6.8cm] (defi) {DeFi：DEX、借贷、稳定币、收益聚合、衍生品协议};
  \node[reg node=teal, below=of defi, minimum width=6.8cm] (l2) {执行与互操作：L2 / Rollup、桥、预言机、MEV 基础设施};
  \node[reg node=orange, below=of l2, minimum width=6.8cm] (l1) {结算层：L1 共识与数据可用性};
  \draw[compliance arrow] (app) -- (defi);
  \draw[compliance arrow] (defi) -- (l2);
  \draw[compliance arrow] (l2) -- (l1);
\end{tikzpicture}
\caption{Web3 协议栈简化分层（示意；各项目实际架构可能跨层）}
\label{fig:web3-stack}
\end{figure}

\section{与本书其他部分的衔接}

技术理解是阅读监管与市场章节的前提：同一「代币」在不同层（支付工具、治理凭证、收益权、抵押品）上的功能组合，会改变法律定性。量化与系统化团队若参与链上策略，还需将\textbf{节点延迟、Gas、回滚、预言机刷新频率}纳入与传统资产不同的执行模型。
